\documentclass[lettersize,journal]{IEEEtran}

\usepackage[hidelinks]{hyperref}
\usepackage{amsmath,amssymb,amsfonts}
\usepackage{cite}
\usepackage{siunitx}

\hyphenation{op-tical net-works semi-conduc-tor IEEE-Xplore exo-skeleton}
\def\BibTeX{{\rm B\kern-.05em{\sc i\kern-.025em b}\kern-.08em
    T\kern-.1667em\lower.7ex\hbox{E}\kern-.125emX}}

\begin{document}

\title{Power-Sign-Matched vs Power-Sign-Mismatched Knee Assistance: A Pilot Study Using a Custom-Designed Knee Exoskeleton}

\author{Juncheng Zhou,
        Weibo Gao,
        and Hao Su$^\dagger$\thanks{All authors are with the Lab of Biomechatronics and Intelligent Robotics, New York University, New York, NY 11201, USA. $^\dagger$Corresponding author: hao.su@nyu.edu.}}

\maketitle

\begin{abstract}
\textbf{Scientific problem:} The closest knee evidence suggests that power-aligned assistance can reduce joint-level demand, but direct causality is still missing. In slope walking, power-inspired knee assistance reduced biological knee power by about 22\% uphill and 20\% downhill, with a 13.4\% uphill metabolic reduction, while group-level EMG reduction remained limited \cite{park2020hingefree}. A context-adaptive uphill controller reported 19\% knee-extensor EMG reduction and 5.8\% metabolic reduction at 8.5$^\circ$ \cite{lee2023context}. However, no healthy-subject knee study has directly compared \emph{power-sign matched} versus \emph{power-sign mismatched} assistance as the primary independent variable under controlled amplitude/work budgets.
\textbf{Scientific question:} does neuromuscular benefit arise from true sign alignment with biological knee power, or from net energy injection despite partial mismatch?
\textbf{Method goal:} We propose a knee-only crossover protocol that explicitly tests this sign effect under identical phase windows and matched delivered-work constraints. The core criterion is
$\operatorname{sign}(P_{\mathrm{exo},k})=\operatorname{sign}(P_{\mathrm{bio},k})$,
where $P_{\mathrm{exo},k}=\tau_{\mathrm{exo},k}\omega_k$ and
$P_{\mathrm{bio},k}=M_{\mathrm{bio},k}\omega_k$.
\textbf{Task objective:} quantify the isolated sign effect on phase-specific quadriceps EMG (RMS and iEMG), then test whether this effect translates to metabolic and biomechanical outcomes. We focus on uphill/downhill tasks to amplify positive and negative knee-power phases and maximize sensitivity to sign-dependent assistance.
\end{abstract}

\begin{IEEEkeywords}
Knee exoskeleton, electromyography, biomechanical power, assistance timing, human--robot interaction.
\end{IEEEkeywords}

\section{Introduction}
Lower-limb exoskeleton assistance is highly sensitive to \emph{when} torque is delivered, not only \emph{how much} torque is delivered. This is well established in hip and ankle studies, where timing optimization can strongly change metabolic outcomes and joint-level mechanics \cite{ding2016timing,galle2017ankle,grimmer2019comparison,lim2019tro}. Knee assistance remains harder to optimize: benefits are often task-dependent, and stronger effects are usually reported in uphill/high-demand conditions than in level walking \cite{park2020hingefree,mclean2019energetics,lee2023context}.

Existing knee evidence already suggests this mechanism gap. Park et al. reported substantial biological power unloading (uphill $\sim$22\%, downhill $\sim$20\%) and uphill metabolic reduction (13.4\%), but without a strong group-level EMG decrease \cite{park2020hingefree}. Lee et al. reported a 19\% knee-extensor EMG reduction and 5.8\% metabolic reduction at 8.5$^\circ$ uphill with context-adaptive assistance \cite{lee2023context}. These studies indicate that demand-matched assistance can work, yet they do not provide a direct matched-versus-mismatched sign comparison under controlled budgets. Therefore, the field still lacks a causal sign-isolation test.

We therefore frame the problem at the power level:
\begin{equation}
P_{\mathrm{exo},k}(t)=\tau_{\mathrm{exo},k}(t)\,\omega_k(t), \qquad
P_{\mathrm{bio},k}(t)=M_{\mathrm{bio},k}(t)\,\omega_k(t),
\label{eq:power_defs_intro}
\end{equation}
where $\tau_{\mathrm{exo},k}$ is exoskeleton knee torque, $M_{\mathrm{bio},k}$ is biological knee moment, and $\omega_k$ is knee angular velocity. Because the sign of power depends on both torque and velocity, torque direction alone cannot determine whether assistance is mechanically cooperative.

This manuscript defines the scientific question and method framework for a dedicated power-sign study in knee exoskeleton control.

\section{Scientific Question and Task Objectives}
\subsection{Scientific Question}
The central question is:
\begin{quote}
\emph{Under matched assistance magnitude and matched phase windows, does power-sign-matched knee assistance reduce phase-specific muscle burden more than power-sign-mismatched assistance?}
\end{quote}

Formally, define the instantaneous sign-consistency indicator
\begin{equation}
\chi(t)=\operatorname{sign}\!\left(P_{\mathrm{exo},k}(t)\,P_{\mathrm{bio},k}(t)\right).
\label{eq:chi}
\end{equation}
Then $\chi(t)=+1$ indicates sign-matched assistance and $\chi(t)=-1$ indicates sign-mismatched interaction.

\subsection{Primary and Secondary Objectives}
Primary objective (EMG): for target muscle set $\mathcal{M}_{\mathrm{quad}}=\{\mathrm{VM},\mathrm{VL},\mathrm{RF}\}$, test whether sign-matched conditions produce lower phase-specific EMG than sign-mismatched conditions.

Secondary objectives:
\begin{enumerate}
    \item metabolic outcome difference between matched and mismatched conditions;
    \item biological knee moment/power/work reduction in target windows;
    \item cross-joint compensation at hip/ankle (biomechanical safety check).
\end{enumerate}

\subsection{Hypotheses}
Let $Y$ denote an outcome and $\Delta Y = Y_{\mathrm{matched}}-Y_{\mathrm{mismatched}}$ under matched window and matched budget constraints.
\begin{equation}
\text{H1: } \Delta \mathrm{EMG}_{\mathrm{phase}} < 0,
\qquad
\text{H2: } \Delta \dot E_{\mathrm{met}} < 0,
\label{eq:hypothesis_h1h2}
\end{equation}
with stronger effects expected in slope tasks where knee power demand is larger.
In parallel, we test whether effects remain after controlling for delivered mechanical budgets:
\begin{equation}
\text{H3: } \Delta \mathrm{EMG}_{\mathrm{phase}}<0 \;\text{under}\; \eqref{eq:match_tau}\text{ and }\eqref{eq:match_power}.
\label{eq:h3}
\end{equation}

\section{Method Framework}
\subsection{Biomechanical Definitions}
For each sample $t$:
\begin{align}
M_{\mathrm{bio},k}(t) &= M_{\mathrm{net},k}(t)-\tau_{\mathrm{exo},k}(t),
\label{eq:mbio}\\
P_{\mathrm{bio},k}(t) &= M_{\mathrm{bio},k}(t)\,\omega_k(t),
\label{eq:pbio}\\
P_{\mathrm{exo},k}(t) &= \tau_{\mathrm{exo},k}(t)\,\omega_k(t).
\label{eq:pexo}
\end{align}
Here $M_{\mathrm{net},k}$ is inverse-dynamics net knee moment.

To avoid sign chatter near zero crossings, define a dead-zone threshold $\delta_P>0$ using baseline cycles:
\begin{equation}
\mathcal{W}^{+}=\{\phi: P_{\mathrm{bio},k}^{0}(\phi)>\delta_P\},
\qquad
\mathcal{W}^{-}=\{\phi: P_{\mathrm{bio},k}^{0}(\phi)<-\delta_P\},
\label{eq:phase_windows}
\end{equation}
where $\phi\in[0,100\%]$ is gait phase and $P_{\mathrm{bio},k}^{0}$ is the baseline biological power profile.

\subsection{Task and Condition Design}
Based on literature showing amplified knee-demand effects on slopes \cite{park2020hingefree,lee2023context,nuckols2020mechanics}, we use two primary tasks:
\begin{enumerate}
    \item uphill walking: dominant positive-power window (typically early stance),
    \item downhill walking: dominant negative-power window (typically early stance).
\end{enumerate}

Default initial windows are set from prior knee slope evidence and then individualized by each subject's baseline zero-crossings:
\begin{equation}
\mathcal{W}_{\mathrm{up}}^{+}\approx [0,30]\%\ \text{GC},\qquad
\mathcal{W}_{\mathrm{down}}^{-}\approx [0,25]\%\ \text{GC}.
\label{eq:default_windows}
\end{equation}

Planned conditions per task:
\begin{enumerate}
    \item no exoskeleton,
    \item exoskeleton unpowered (zero-work),
    \item matched-low, matched-medium,
    \item mismatched-low, mismatched-medium.
\end{enumerate}
An optional condition can be added for timing disentanglement: same-sign time-shifted assistance.

\subsection{Sample Size and Data-Acquisition Plan}
The first full study targets 20--24 completed healthy participants, with recruitment expanded if needed to account for drop-out or unusable EMG/metabolic trials. This target is chosen because existing knee studies are often small and because the planned interaction terms (condition $\times$ task $\times$ muscle) require higher within-subject statistical power \cite{park2020hingefree,mclean2019energetics,lee2023context}.

Recommended instrumentation and rates are: motion capture (100--200 Hz), force plates or instrumented treadmill (1000--2000 Hz), surface EMG (1000--2000 Hz), and synchronized exoskeleton sensing/control logs ($\ge$500 Hz). Each condition uses a 6-min block with 3--4 min adaptation and a late steady segment for analysis.

\subsection{Sign-Matched vs Sign-Mismatched Control Law}
Let $g(\phi)\ge 0$ be a fixed envelope and $A\in\{A_{\ell},A_m\}$ be amplitude level. Define reference sign from baseline biological power:
\begin{equation}
s_{\mathrm{bio}}(\phi)=\operatorname{sign}\!\left(P_{\mathrm{bio},k}^{0}(\phi)\right).
\label{eq:sbio}
\end{equation}
Assistance command is
\begin{equation}
\tau_{\mathrm{cmd}}(\phi;A,m)=A\,g(\phi)\,s_{\mathrm{bio}}(\phi)\,m,
\qquad m\in\{+1,-1\},
\label{eq:tau_cmd}
\end{equation}
where $m=+1$ is sign-matched and $m=-1$ is sign-mismatched. This flips sign while preserving waveform magnitude and timing envelope.

\subsection{Budget-Matching Constraints}
To isolate sign effect, matched and mismatched conditions are constrained by delivered-budget matching in each target window $\mathcal{W}_j$:
\begin{align}
\left|\int_{\mathcal{W}_j}\!|\tau^{(+)}_{\mathrm{exo},k}(t)|dt-
\int_{\mathcal{W}_j}\!|\tau^{(-)}_{\mathrm{exo},k}(t)|dt\right| &\le \epsilon_{\tau},
\label{eq:match_tau}\\
\left|\int_{\mathcal{W}_j}\!|P^{(+)}_{\mathrm{exo},k}(t)|dt-
\int_{\mathcal{W}_j}\!|P^{(-)}_{\mathrm{exo},k}(t)|dt\right| &\le \epsilon_{P}.
\label{eq:match_power}
\end{align}
Superscripts $(+)$ and $(-)$ denote matched and mismatched cases.

\subsection{Outcome Metrics}
For muscle $r$ in phase window $\mathcal{W}_j$ with duration $T_j$:
\begin{align}
\mathrm{RMS}_{r,j} &=
\sqrt{\frac{1}{T_j}\int_{\mathcal{W}_j} e_r^2(t)\,dt},
\label{eq:rms}\\
\mathrm{iEMG}_{r,j} &=
\int_{\mathcal{W}_j} |e_r(t)|\,dt,
\label{eq:iemg}
\end{align}
where $e_r(t)$ is processed EMG envelope.

To directly quantify the sign intervention in each condition, we compute
\begin{equation}
R_{\mathrm{sign},j}=
\frac{1}{T_j}\int_{\mathcal{W}_j}\mathbf{1}\!\left[\chi(t)=+1\right]dt,
\label{eq:rsign}
\end{equation}
which is the matched-sign occupancy in the target window.

Additional mechanics metrics:
\begin{align}
W_{\mathrm{bio},j} &= \int_{\mathcal{W}_j} P_{\mathrm{bio},k}(t)\,dt,
\qquad
W_{\mathrm{exo},j}=\int_{\mathcal{W}_j} P_{\mathrm{exo},k}(t)\,dt,
\label{eq:work_terms}\\
R_{\mathrm{align},j} &=
\frac{1}{T_j}\int_{\mathcal{W}_j}
\mathbf{1}\!\left[P_{\mathrm{exo},k}(t)P_{\mathrm{bio},k}(t)>0\right]dt.
\label{eq:align_ratio}
\end{align}

\subsection{Statistical Model}
Primary EMG analysis uses a linear mixed-effects model:
\begin{equation}
\log(\mathrm{RMS}_{r,j}) =
\beta_0 + \beta_1 I_{\mathrm{match}} + \beta_2 I_{\mathrm{mag}} + \beta_3 I_{\mathrm{task}} +
\beta_4(I_{\mathrm{match}}\times I_{\mathrm{task}}) + b_{0,s} + b_{1,s}I_{\mathrm{match}} + \varepsilon,
\label{eq:lmm}
\end{equation}
where $s$ indexes subject, $b_{0,s},b_{1,s}$ are random effects, and $\varepsilon$ is residual error. The same structure is applied to iEMG and key secondary outcomes. Planned contrasts prioritize matched vs mismatched within each task (uphill, downhill).

\section{Task Goal for the Full Study}
The full study is designed to deliver a direct, causal answer to the current literature gap:
\begin{enumerate}
    \item determine whether power-sign matching itself reduces phase-specific knee-related EMG burden;
    \item determine whether this neuromuscular effect translates into metabolic and biomechanical benefit;
    \item provide a reproducible sign-explicit protocol for future knee exoskeleton controller comparisons.
\end{enumerate}

\bibliographystyle{IEEEtran}
\bibliography{Reference}

\end{document}
